% MAIN SNAPSHOT — LSPV/RMV::1056::2026-10-02T10:56:00-07:00
\documentclass[11pt]{article}
\usepackage[margin=.86in]{geometry}
\usepackage{amsmath,amssymb}
\begin{document}
\section*{LSPV-5 MAIN Snapshot 1056}

\textbf{STATUS: OPEN.} P versus NP remains open. No unrestricted SAT lower bound,
arbitrary-algorithm acquisition theorem, nonrelativizing hinge, or nonalgebrizing hinge is established.

\section*{CORPUS RECONCILIATION}
\textbf{Corpus surveyed.}
The full current repository tree was enumerated before substantive work. This run reviewed root
\texttt{main.tex}; AGENTS/CODEX/pnp-research governance; TASK-001--007; verified MAIN history through
1012; materially relevant LSPV-1--4 reports including MATRIXCERT1, TURANEX1, RAMCOMP1 and predecessor
TRACECLOSE1; BITRANK/BILINREC/DIGCORR/SPARSEEX; PRIME-LQP/PRIMECRT/PH-PROV-CLOSE;
HEC/VSC-HEC/G-PPC-HEC/CRS-PPC/CCSPC/RCCN/Search--G; ACR1, SBACR1 and EVACR1; current
GROK/CLAUDE and independent-review artifacts; governance; and cited ACR/product-span experiments and
raw results. The local preflight executable could not run because the active container has no
repository checkout; GitHub-side governance and readback checks were performed directly.

\textbf{New local theorem, independent red-team review pending.}
STRATFLAT1 proves a source-strategy flattening lemma. For polynomially many alternating rounds
\[
\exists R\;\forall c_1\exists o_1\cdots\forall c_T\exists o_T\;V,
\]
if the existential openings are supplied by polynomial-time Skolem functions
\[
o_i=F_i(x,R,c_{\le i}),
\]
then the game is equivalent to
\[
\boxed{\exists R\;\forall C\;W(x,R,C)}
\]
with \(W\in P\), hence lies in \(\Sigma_2^P\). More generally, if a winning strategy has a
polynomial-size description \(S\) and polynomial-time evaluator, existence of such a strategy is
also a \(\Sigma_2^P\) relation:
\[
\exists R,S\;\forall C\;V(x,R,C,\operatorname{Eval}(S,C)).
\]
Under \(P=NP\), PH-PROV-CLOSE1 therefore closes these succinct-strategy cases.

\textbf{Source-matched consequence.}
SBACR1's canonical source response
\[
F(w,s)=\operatorname{RREFBasis}(S_s(w))
\]
is polynomial-time computable from the static Boolean-grid source by interpolation, orbit
evaluation, binary labeling and Gaussian elimination. Thus adding polynomially many source-bound
slice challenges does not create a genuinely PSPACE-like resource. EVACR1 independently shows that
ordinary evolving restrictions admit a succinct compositional history quotient.

\textbf{Computational control.}
STRATFLAT1 exhaustively checked the source-bound game for
\[
(p,h,m,T)=(17,4,4,4)
\]
on 917,504 source/challenge paths and for
\[
(29,5,5,3)
\]
on 658,560 paths. Interactive canonical openings always equaled flattened direct evaluation.
Local response nonuniqueness was present on 220/224 and 839/840 source-slices respectively, because
rank-at-least-two slices admit multiple ordered basis encodings. Canonical RREF still supplied a
single polynomial-time winning selector. Status: \textbf{COMPUTATIONAL} finite control.

\textbf{Conflict resolved.}
The 1012 suggestion that polynomial alternation itself might be the first surviving provenance zone
is too weak. Polynomial alternation only matters if every winning strategy lacks a polynomial-size
or polynomial-time-evaluable representation. Otherwise the strategy can be Skolemized/flattened
back into fixed PH.

\textbf{Stale claims removed/downgraded.}
Generic authenticated replay, response-signature diversity, explicit adaptive-tree size, ordinary
evolving restrictions, and polynomial alternation with a succinct strategy are now all rejected as
generic acquisition mechanisms. They remain useful instrumentation and hostile controls.

\textbf{Falsifications retained.}
REACQ1, TRACECLOSE1, PH-PROV-CLOSE1, ACRTREE1, SBACR1 reread/recompute, EVACR1 compositional
collapse, COVERBREAK1, EXCESS1, BITRANK1's characteristic-transfer failure, SGLAUD1, ordinary
switching as universal transfer, generic canonical binding, raw incidence/tape geometry, and prime
magnitude/factorization remain retained.

\textbf{Adopted/rejected handoffs.}
Adopt STRATFLAT1 as a local closure theorem pending independent review, and retain QBF strategy
extraction/synthesis literature as exact context for succinct strategies. Reject QBF/DQBF proof
lower bounds or PSPACE/NEXP strategy problems as direct SAT lower-bound transfers absent a
model-preserving compiler.

\section*{Literature reconciliation}
Chandra--Kozen--Stockmeyer give \(AP=PSPACE\), establishing the genuine polynomial-alternation
boundary. Beyersdorff--Blinkhorn--Mahajan show QBF strategies can be represented succinctly inside
proofs. Heule--Seidl--Biere extract Skolem functions from QRAT proofs. Akshay et al. (SAT 2026)
study proof systems permitting efficient extraction of Skolem/Herbrand functions and game values.
Santhanam--Williams relate faster QBF algorithms to circuit lower bounds. Together these results
support the present diagnosis: the relevant resource is strategy representation/evaluation
complexity, not quantifier count alone.

\section*{Shortest surviving proof chain}
\[
\boxed{\text{exact source-typed SU/vMS family}}
\to
\boxed{\operatorname{rank}_{\mathbb F_2}\Phi=\omega(m)\text{ or classified exceptions}}
\to
\boxed{\text{forced SAME\_REALIZATION acquisition not encoded by a succinct source/strategy}}
\]
\[
\to
\boxed{\text{charged reconstruction}}
\to
\boxed{\text{nonrelativizing/nonalgebrizing hinge}}
\to
\boxed{\text{unconditional contradiction}}.
\]
Every load-bearing boxed arrow remains OPEN or CONDITIONAL.

\section*{Structural branch}
BILINREC1, DIGCORR1 and SPARSEEX1 retain their exact declared scopes. BITRANK1 still requires an
intrinsically characteristic-two certificate rather than complex Gram rank. MATRIXCERT1 supplies
legal-slice pair-incidence certificates. EXCESS1/EXCESSDUAL1 subtract the universal local-nullity
baseline, and TURANEX1 converts universal low excess into generalized-weight support constraints.
COVERBREAK1 remains the mandatory finite survivor. None of these statements forces arbitrary SAT
runtime.

\section*{Causal/provenance branch}
The closure ladder is now stronger:
\[
\text{semantic quotient}
\to
\text{deterministic replay}
\to
\text{fixed-PH replacement}
\to
\text{succinct-strategy flattening}.
\]
A source-bound replay candidate survives only if it proves that every legal winning strategy has
superpolynomial representation/evaluation complexity and that this exact game is forced on every
hypothetical static SAT decider. No such theorem is known. This is at least as difficult as a new
strategy/circuit lower bound plus an arbitrary-algorithm compiler.

\section*{ONE SINGLE TEAM TARGET}
Return the primary effort to one exact reconstruction-compatible SU/vMS polynomial-bit orbit.
Construct the source-transcribed binary operator \(\Phi\) and legal slice matrices, and prove or
falsify
\[
\boxed{\operatorname{rank}_{\mathbb F_2}\Phi=\omega(m)}
\]
using odd minors, binary dual annihilators, MATRIXCERT1, EXCESS/TURANEX1 generalized weights and
COVERBREAK1 in one pipeline.

As a mandatory hostile side test on the same family, explicitly search for compact response
representations: algebraic evaluators, Skolem circuits, decision diagrams, or source/history
quotients computing all legal replay responses. If any polynomial representation survives,
provenance remains instrumentation only.

\section*{What this does not prove}
STRATFLAT1 is a closure theorem, not a separation. The finite source tests do not prove asymptotics.
A strategy-representation lower bound would still need an arbitrary-static-SAT necessity theorem,
and the nonrelativizing/nonalgebrizing terminal step remains OPEN. P versus NP remains OPEN.

\begin{thebibliography}{9}
\bibitem{CKS81} A. K. Chandra, D. C. Kozen, L. J. Stockmeyer,
Alternation, J. ACM 28(1), 1981, 114--133.
\bibitem{BBM21} O. Beyersdorff, J. Blinkhorn, M. Mahajan,
Building Strategies into QBF Proofs, J. Automated Reasoning 65, 2021.
\bibitem{HSB14} M. Heule, M. Seidl, A. Biere,
Efficient Extraction of Skolem Functions from QRAT Proofs, FMCAD 2014.
\bibitem{Akshay26} S. Akshay et al.,
Proof Systems for QBF Synthesis: Extracting Skolem and Herbrand Functions, SAT 2026.
\bibitem{SW15} R. Santhanam, R. Williams,
Beating Exhaustive Search for Quantified Boolean Formulas and Connections to Circuit Complexity,
SODA 2015.
\end{thebibliography}

\end{document}
% END_OF_VERIFIED_MAIN::1056
